\section{Conclusion}
\label{sec:conclusion}

\textbf{Problem and motivation:} existing LLM-for-CPS security studies are offline or simulation-only, leaving open whether language models can conduct effective closed-loop security analysis of live industrial systems---both as red team (attacker) and blue team (defender).

\textbf{Insight:} CPS security analysis benefits from a closed-loop red team / blue team framework where the same LLM serves both roles against a live controller. This requires an online feedback loop, controlled context ablation, configurable defense chains, and cross-model comparison---evaluated on real PLC hardware.

\textbf{Solution:} CPSForge is a general-purpose PLC-in-the-loop framework for automated LLM-driven red/blue team analysis of cyber-physical systems. The framework places a local LLM in the control loop as both red team (online MITM attacker) and blue team (anomaly defender) using the same model and context pipeline. All writes pass through a mandatory safety shield and configurable defense chain. The framework is configuration-driven: any CPS connected to a supported PLC can be analyzed by providing a scene configuration file.

We evaluate CPSForge across three Factory~I/O scenes on a real Siemens S7 PLC, organized into six research questions:
\begin{itemize}[leftmargin=1.5em, nosep]
    \item \textbf{RQ1 (Context):} How does operational context affect red team effectiveness?
    \item \textbf{RQ2 (Fine-Tuning):} How does task-specific training change both red and blue team capability?
    \item \textbf{RQ3 (Feedback Loop):} Is closed-loop feedback essential for LLM red teaming?
    \item \textbf{RQ4 (Blue Team):} Which blue team configurations are most effective?
    \item \textbf{RQ5 (Cross-Model):} Do findings generalize across models?
    \item \textbf{RQ6 (Frontier Model):} Does a frontier API model exhibit qualitatively different attack behavior?
\end{itemize}

Key findings: (1)~the online MITM achieves up to 54\% ASR (3B local) and 100\% ASR (GPT-4o-mini) on Sorting by Weight---lower bounds under the mandatory safety shield---while the static one-shot design generates zero valid proposals; (2)~the LLM blue team blocks 100\% of attack proposals on Sorting by Height, including cross-model blocking of frontier-model attacks, but with 95--100\% false-positive rate on legitimate writes (226/229 probes; Section~\ref{sec:discussion}); (3)~red team ASR scales with model size (0\% at 1.7B, 14\% at 3B, 100\% for GPT-4o-mini on Sort.~Weight minimal) but the relationship is scene-dependent---on Sort.~Height, the 3B local model outperforms the frontier model (13.3\% vs.\ 0\%); (4)~QLoRA fine-tuning improves red team ASR from 0\% to 37\% on Level Control while preserving 100\% LLM defender blocking; (5)~the self-deterrence effect under full operational context, observed with 3B models, vanishes entirely at frontier model scale (GPT-4o-mini generates 10/10 proposals in every cell); (6)~context depth effects are scene-dependent: full context suppresses attacks on continuous-control scenes (3B models only) but enables them on discrete state-machine scenes.

Seven contributions: (C1)~a closed-loop LLM red team model, (C2)~a context-ablation framework, (C3)~a before/after fine-tuning study on both red and blue team capability, (C4)~an LLM blue team agent with rule-based and combined configurations, (C5)~a cross-model robustness evaluation spanning 1.7B to frontier scale, (C6)~a frontier model scaling study demonstrating self-deterrence elimination and scene-dependent ASR inversion, and (C7)~a comprehensive PLC-in-the-loop evaluation (333~runs). Additionally, Tier~3 live logic MITM attacks demonstrate that a 3B-parameter LLM can generate minimal SCL code modifications (comparison inversions) that compile without errors and cause severe physical disruption on all three evaluation scenes---validating the end-to-end pipeline from LLM code analysis to real PLC deployment.

Within the scope of this evaluation---one PLC family (Siemens S7-1200) and three software-emulated scenes (15--35 tags)---CPSForge is, to our knowledge, the first closed-loop LLM red/blue team framework for CPS security analysis evaluated on a real PLC. Extending to other PLC families (Allen-Bradley, Schneider), larger-scale plants (50+ tags with cross-subsystem dependencies), and physical testbeds (SWaT-class infrastructure) is the primary direction for future work. The framework, configurations, and all experiment artifacts are released for reproducibility and extension to other industrial systems.
